%% =====================================================================
%%  proyecto-de-ley.tex  --  Plantilla de PROYECTO DE LEY (MDD)
%% ---------------------------------------------------------------------
%%  La estructura la fija el esquema `proyecto-de-ley` del ecosistema
%%  (Reglamento del Congreso art. 75; Manual de Tecnica Legislativa 2021
%%  sec. X y VI) y el compilador la exige: cada pieza se abre con
%%  \componente{clave}. Metodo: prompts/06 legislativo/prompt_08.
%%  Las cifras de la exposicion de motivos llevan fuente, periodo y unidad;
%%  ninguna norma se cita sin verificarla sobre el texto (prompts/docs/dominios/legislativo.md §3).
%% =====================================================================
\documentclass{proyecto-de-ley}
%% =====================================================================
%%  metadata.tex  --  DATOS del proyecto de ley
%% =====================================================================
\setmeta{numero}{001}
\setmeta{anio}{2026}
% \setmeta{fecha}{11 de setiembre del 2026}   % comentado: usa la fecha de hoy
\setmeta{firmante}{\meta{autor}}   % por defecto firma el titular del perfil
\setmeta{cargofirmante}{Diputado de la Rep\'ublica}
\setmeta{grupo}{Grupo Parlamentario [Inserta aqu\'i]}

\esquemacondicion{modificatoria}   % la propuesta modifica una norma vigente: cuadro comparativo
% \esquemacondicion{ambiental}     % solo si la materia tiene incidencia ambiental
% \esquemacondicion{agenda}        % relacion con la Agenda Legislativa / Acuerdo Nacional
% \esquemacondicion{anexo}

\begin{document}

%% ── Presentacion (Manual sec. X.c.1) ──────────────────────────────────
\componente{presentacion}
\noindent El \meta{grupo}, a iniciativa del congresista de la Rep\'ublica que
suscribe, \meta{firmante}, en ejercicio del derecho de iniciativa legislativa
que le confiere el art\'iculo 107 de la Constituci\'on Pol\'itica del Per\'u, y de
conformidad con los art\'iculos 75 y 76 del Reglamento del Congreso de la
Rep\'ublica, propone el siguiente:

%% ── Titulo ────────────────────────────────────────────────────────────
\componente{titulo}
\titulodelaley{Proyecto de ley que [Inserta aqu\'i: objeto, en una frase]}

%% ── Formula legal (Manual sec. VI) ────────────────────────────────────
\componente{formula}
\begin{center}{\seriesemi EL CONGRESO DE LA REP\'UBLICA\\ HA DADO LA LEY SIGUIENTE:}\end{center}

\componente{sustantiva}
\articulo{1}{Objeto de la ley}
[Inserta aqu\'i: una regla por art\'iculo; sin fundamentaci\'on dentro del articulado.]

\articulo{2}{Finalidad}
[Inserta aqu\'i.]

\componente{dcf}
\bloquedisposiciones{finales}
\disposicion{Primera}{Vigencia}
La presente ley entra en vigencia al d\'ia siguiente de su publicaci\'on en el
diario oficial \emph{El Peruano}.

\componente{dcm}
\bloquedisposiciones{modificatorias}
\disposicion{\'Unica}{Modificaci\'on del art\'iculo [n] de la Ley N.° [Inserta aqu\'i]}
Modif\'icase el art\'iculo [n] de la Ley N.° [Inserta aqu\'i], en los t\'erminos
siguientes: [texto \'integro del art\'iculo tal como queda].

\componente{firmas}
\firmaslegislativas

%% ── Exposicion de motivos (art. 75; Manual sec. X.c.4) ────────────────
\componente{exposicion}

\componente{fundamentos}
[El problema con su evidencia; el estado actual de la situaci\'on f\'actica o
jur\'idica; el nuevo estado que genera la propuesta; necesidad, viabilidad y
oportunidad; el marco normativo; las opiniones, cuando corresponda.]

\componente{antecedentes}
[Cada antecedente parlamentario por su n\'umero de expediente y estado,
verificado en la fuente oficial.]

\componente{efecto}
[Qu\'e art\'iculos o partes de art\'iculos se modifican o derogan, con precisi\'on.]

\componente{costo-beneficio}
[Sectores beneficiados y perjudicados; efectos monetarios y no monetarios;
impacto econ\'omico; impacto presupuestal cuando corresponda. Lo que no se
puede cuantificar se declara como no cuantificado.]

%% ── Cuadro comparativo (condicion `modificatoria`) ────────────────────
\componente{comparativo}
\begin{longtable}{@{}p{0.48\textwidth}p{0.48\textwidth}@{}}
  \toprule
  {\seriesemi Texto vigente} & {\seriesemi Texto propuesto} \\
  \midrule
  \endhead
  Art\'iculo [n].\ [Inserta aqu\'i] & Art\'iculo [n].\ [Inserta aqu\'i] \\
  \bottomrule
\end{longtable}

\end{document}
